\documentclass[a4paper,11pt]{article}
\usepackage[ngerman]{babel}
\usepackage[top=3cm,bottom=3cm,left=3cm,right=3cm]{geometry} %NICHT IN KEMIE2


%\usepackage{microtype} %schöneres typesetting  mit [stretch=10] für nur 1% anstelle von 2%
%\usepackage{lmodern} %Latin Modern FONT
%\usepackage{pifont} %schönere Font



\usepackage{amsmath}
\usepackage{nicefrac}
\usepackage{amssymb}
%\usepackage{mathtools}

\usepackage{titlesec} %Um Section, etc bearbeiten zu können
\usepackage{setspace}
\usepackage{enumitem} %enumerate
\usepackage{multicol}
\usepackage[many]{tcolorbox}
\usepackage{hyperref}
\usepackage{arydshln}

%\renewcommand{\ttdefault}{lmtt} %Schriftart wird geändert zu Latin Modern Typewriter


\hypersetup{hidelinks,
			colorlinks=true,
			linkcolor=black,
			citecolor=miudiutex,
			urlcolor=miugray2
			}

\setlength{\parindent}{0pt}
\setcounter{section}{0}


\titlespacing*{\section}{0pt}{20mm}{6mm}
\titlespacing*{\subsection}{0pt}{6mm}{3mm}
\titlespacing*{\subsubsection}{0pt}{3mm}{0mm}



%\definecolor{miudiutex}{HTML}{2f3d39}
\definecolor{miudiutex}{HTML}{1d2926}
\definecolor{miugray}{HTML}{b4cccf}
\definecolor{miugray2}{HTML}{4d7365}
\definecolor{red}{HTML}{cf1f5c}
\definecolor{green}{HTML}{00A36C}
\definecolor{delphinidin}{HTML}{315794}
\definecolor{pelargonidin}{HTML}{b33807}
\definecolor{cyanidin}{HTML}{ba0746}
\definecolor{petunidin}{HTML}{b56fed}



\raggedcolumns %wichtig für Multicolumns


%================================
% SYMBOLS
%================================

\newcommand{\RA}{$\rightarrow$~}
\newcommand{\RL}{$\rightleftharpoons$~}
\newcommand{\HARP}{\ding{226}~}
%$\xrightarrow{2}$ %Pfeil mit 2 drüber


%================================
% SHORT COMMANDS (BOXES ; ETC)
%================================


\newcommand{\notiz}[1]{\begin{notizz}{}{}\begin{center}
			{\color{miugray2!50!miudiutex}#1} \end{center}\end{notizz}}
\newcommand{\zit}[2]{\begin{zitat*}{#1}{}\small #2 \end{zitat*}}

\newcommand{\frage}[2]{\begin{qst*}{#1}{}\small #2 \end{qst*}}

\newcommand{\unten}{\end{center} \tcblower \begin{center}}


\newcommand*\circled[1]{\tikz[baseline=(char.base)]{
            \node[shape=circle,draw,inner sep=0.5pt,miudiutex] (char) {\tiny #1};}}
            
\newcommand{\miusquare}{{\color{miugray2!50}\scriptsize $\blacksquare$}~}
\newcommand{\miusquarp}{{\color{petunidin!50}\scriptsize $\blacktriangleright$}~}
\newcommand{\niu}{\item[\miusquare]}
\newcommand{\nip}{\item[\miusquarp]}

%%%%% ABSAETZE

\newcommand{\pps}{\par \vspace*{1mm}}
\newcommand{\pp}{\par \vspace*{3mm}}
\newcommand{\ppbig}{\par \vspace*{8mm}}
\newcommand{\pphuge}{\par \vspace*{10mm}}

%================================
% FRAGE BOX
%================================

\tcbuselibrary{theorems,skins,hooks}
\newtcbtheorem{qst}{Frage:}
{%
	enhanced,
	breakable,
	colback = gray!15!white,
	frame hidden,
	boxrule = 0sp,
	borderline west = {2.5pt}{0pt}{red!70},
	sharp corners,
	detach title,
	before upper = \tcbtitle\par\smallskip,
	coltitle = black,
	fonttitle = \bfseries\sffamily,
	description font = \mdseries,
	separator sign none,
	segmentation style={solid, gray!50!black},
}
{th}




%================================
% ZITAT BOX
%================================

\tcbuselibrary{theorems,skins,hooks}
\newtcbtheorem{zitat}{Zitat}
{%
	enhanced,
	breakable,
	colback = gray!15!white,
	frame hidden,
	boxrule = 0sp,
	borderline west = {2.5pt}{0pt}{black!70},
	sharp corners,
	detach title,
%	before upper = \tcbtitle\par\smallskip,
	coltitle = gray!50!black,
	fonttitle = \bfseries\sffamily,
	description font = \mdseries,
%	separator sign none,
	segmentation style={solid, gray!50!black},
}
{th}



%================================
% SIMPLE SCHWARZER RAHMEN BOX
%================================


\newcommand{\boxfr}[1]{
	\begin{tcolorbox}[
	drop shadow southeast,
	enhanced,
	colframe=black!50,
	colback=white,
	boxrule = 1pt, 
%	toprule=0pt, bottomrule=0pt,rightrule=0pt,leftrule=0pt,
	boxsep=5pt,
	arc=1pt,
	]
	\begin{center}
	#1
	\end{center}
	\end{tcolorbox}
}

\definecolor{miudiutex}{HTML}{1d2926}
\definecolor{miugray}{HTML}{b4cccf}
\definecolor{miugray2}{HTML}{4d7365}
\definecolor{white2}{HTML}{f5f7f8}
\definecolor{red}{HTML}{cf1f5c}
\definecolor{green}{HTML}{00A36C}
\definecolor{delphinidin}{HTML}{315794}
\definecolor{uniblau}{HTML}{004e9f}
\definecolor{unigelb}{HTML}{fcba00}
\definecolor{unigrau}{HTML}{909085}
\definecolor{pelargonidin}{HTML}{b33807}
\definecolor{cyanidin}{HTML}{ba0746}
\definecolor{paeonidin}{HTML}{d15ca6}
\definecolor{petunidin}{HTML}{b56fed}
\definecolor{malvidin}{HTML}{8a2d8c}

\newcommand{\metermilli}{\mskip3mu \text{mm}}
\newcommand{\cm}{\mskip3mu \text{cm}}
\newcommand{\m}{ \text{m}}
\newcommand{\lite}{ \text{l}}
\newcommand{\kg}{ \text{kg}}
\newcommand{\g}{ \text{g}}
\newcommand{\km}{ \text{km}}
\newcommand{\h}{ \text{h}}
\newcommand{\s}{ \text{s}}
\newcommand{\tonne}{ \text{t}}
\usepackage[utf8]{inputenc}
\usepackage[T1]{fontenc}
\usepackage{lmodern}

\usepackage[
    activate={true,nocompatibility},
    final,
    tracking=true,
    kerning=true,
    spacing=true,
    factor=1100,
    stretch=30,
    shrink=20
]{microtype}
\usepackage{pdfpages}
\usepackage{awesomebox}
\usepackage{marginnote}
\tcbuselibrary{documentation}
\usepackage{sidenotes}
\usepackage{import}
\usepackage{xifthen}
\usepackage{transparent}
\usepackage[table]{xcolor}


\newcommand{\abbicon}{%
  \protect\tikz[baseline=0.2mm,scale=0.8]{%
    % Das blaue Quadrat mit abgerundeten Ecken
    \fill[uniblau, rounded corners=1.2pt] (0,0) rectangle (0.8em, 0.8em);
    % Der weiße Kreis in der Mitte
    \fill[white] (0.4em, 0.4em) circle (0.12em);
  }%
} 
\usepackage[labelfont={bf},font={color=miudiutex,small},figurename={\abbicon $~$Abbildung}]{caption}


\newcommand{\bckgrnd}{
\AddToHookNext{shipout/background}{
\begin{tikzpicture}[remember picture, overlay]
 \draw[draw=none,fill=uniblau!70, shift={(current page.north west)}] (-3,0) rectangle (23,-3.75);
 \draw[draw=none,fill=miugray, shift={(current page.north east)}] (-7,0) rectangle (3,-3.75);
\end{tikzpicture}
}
}



\newcommand{\incfig}[2]{%
\def\svgwidth{#2\columnwidth}
\import{C://LaTeX-Files/pngs/}{#1.pdf_tex}
}

\renewcommand{\textbf}[1]{\textsf{{\bfseries {#1}}}}



\titleformat{\section}[hang]
		{\LARGE \color{uniblau!75!black} \sffamily \bfseries}
		{\thesection}
		{6mm}
		{}
		[]
\titleformat{\subsection}[hang]
		{\large \sffamily \bfseries \color{black}}
		{\thesubsection}
		{4mm}
		{{\color{miugray}$\blacksquare ~$}}

\titleformat{\subsubsection}[block]
		{\normalsize \bfseries \sffamily \color{miudiutex}}%
		{\normalsize \color{black} \thesubsubsection}
		{2mm}
		{{\color{miugray}$\blacksquare ~$}}
		[ \color{black}]
		
\titleformat{\paragraph}
		{\normalsize \bfseries \sffamily \color{black}}
		{\footnotesize \theparagraph}
		{1mm}
		{}

\pagestyle{empty}
\titlespacing*{\section}{0pt}{20mm}{12mm}
\titlespacing{\section}{0pt}{20mm}{12mm}

\titlespacing*{\subsection}{0pt}{14mm}{4mm}
\titlespacing{\subsection}{0pt}{14mm}{4mm}

\titlespacing*{\subsubsection}{0pt}{6mm}{3mm}
\titlespacing{\subsubsection}{0pt}{6mm}{3mm}

\titlespacing*{\paragraph}{0pt}{6mm}{2mm}
\titlespacing{\paragraph}{0pt}{6mm}{2mm}


\let\oldmarginfigure\marginfigure
\let\endoldmarginfigure\endmarginfigure

\let\oldmarginfigure\marginfigure
\let\endoldmarginfigure\endmarginfigure

\renewenvironment{marginfigure}[1][0pt] % [1] steht für 1 Argument, [0pt] ist der Standardwert
  {\oldmarginfigure[#1]\captionsetup{labelfont={sf,bf,color=uniblau!75!black},font={color=miudiutex,small}}}
  {\endoldmarginfigure}
  
  
  
  
\renewcommand{\labelitemi}{\color{uniblau!70}$\bullet$}
\renewcommand{\labelitemii}{\color{uniblau!70}$\cdot$}
\newcommand{\plmtsquare}{{\color{uniblau!70}$\blacksquare~$}}
\newcommand{\splmt}[1]{{\color{uniblau!75!black} \sffamily #1}}
\newcommand{\fplmt}[1]{{\color{uniblau!75!black} \sffamily \bfseries #1}}

\newcommand{\antwort}[1]{
\begin{tcolorbox}[
	enhanced,
	enforce breakable,
	standard jigsaw,
	boxrule=0.7pt,
	colframe=black,
	sharp corners,
	boxsep=5pt,
	title=\bfseries \sffamily Antwort,
	opacityback=0,
	coltitle=miudiutex,
	attach boxed title to top text left={yshift=-\tcboxedtitleheight/2},
	boxed title style={colback=white,colframe=white},
	bottomrule at break=0pt,
	toprule at break=0pt,
	pad at break=16mm,
]
	\begin{center}
	#1
	\end{center}
	\end{tcolorbox}
}




\rowcolors{2}{miugray!50}{miugray!50}
\arrayrulecolor{white} % Weiße Gitterlinien
\setlength{\arrayrulewidth}{1.5pt} % Etwas dickere Linien
\renewcommand{\arraystretch}{1.4} % Mehr Zeilenhöhe für bessere Lesbarkeit}


%\setlist[itemize]{itemsep=0mm}
\setlist[enumerate]{itemsep=0mm}
\usepackage[table]{xcolor}
\usepackage{pgfplots}
    \pgfplotsset{
        every axis post/.style={
	yticklabel=\empty,
	xticklabel=\empty,
	y label style={at={(axis description cs:-0.2,1.1)}},
	x label style={at={(axis description cs:1.05,0.2)}},
	ticks=none,
	width=6cm,
	axis lines=middle,
%	xlabel near ticks,
        },
    }
\usetikzlibrary{positioning, arrows.meta, calc}

\usepackage{relsize}
\usepackage{cancel}

\newif\ifshowme
%\showmetrue



\begin{document}

%\newgeometry{top=3cm,bottom=3cm,left=4.5cm,right=4.5cm}
\newgeometry{top=2cm,bottom=2cm,left=4cm,right=4cm}


\begin{center}
\LARGE
\color{uniblau}
\textbf{Übung 05} \par \vspace*{2mm}
\large \color{miugray!50!black} \textbf{Prozessbezogene Lebensmitteltechnologie}\par 
\normalsize Zusatzaufgaben
\end{center}
\par 
%\vspace*{6mm}


\color{miudiutex}

%\begin{multicols}{2}
\subsection*{Aufgabe 1}
Welcher (statische) Druck herrscht in der Tiefe h unterhalb der Wasseroberfläche?
\begin{figure}[h]
\centering
\begin{tikzpicture}[>=stealth, font=\sffamily, scale=1.2]

    % Leichte blaue Füllung zur Andeutung des Fluids (modellhaft)
    \fill[miugray!25] (-4, 0) rectangle (4, 4);

    % Wasseroberfläche
    \draw[thick, uniblau] (-4, 4) -- (4, 4);
    
    % Symbol für freie Oberfläche (Wasserspiegel)
    \draw[uniblau, thick] (2.5, 4) -- (2.65, 4.25) -- (2.35, 4.25) -- cycle;
    \draw[uniblau, thick] (2.25, 4.25) -- (2.75, 4.25);
    \draw[uniblau, thick] (2.35, 4.35) -- (2.65, 4.35);

    % Gestrichelte Linien für die Höhenniveaus
    \draw[dashed, thick] (-1.5, 4) -- (1, 4);
    \draw[dashed, thick] (-1.5, 0) -- (1, 0);

    % Pfeil für den Höhenunterschied (h)
    \draw[<->, thick] (0, 0) -- (0, 4);

    % Stromlinie (geschwungener Pfad)
    \draw[->, thick, unigelb] (0, 4) 
        .. controls (2.5, 3) and (1.5, 2) .. (0, 1.5) 
        .. controls (1,1) and (1.5,0.5) .. (0,0)
        node[pos=0.55, right=3mm] {Stromlinie};

    % Beschriftungen Oben (Niveau 1)
    \node[left=2mm] at (-1.5, 4.3) {$h_1 = h$};
    \node[right=2mm] at (0, 4.3) {$p_1 = p_u$};

    % Beschriftungen Unten (Niveau 2)
    \node[left=2mm] at (-1.5, 0) {$h_2 = 0$};
    \node[right=2mm] at (1.5, 0) {$p_2 = ?$};

\end{tikzpicture}

\end{figure}
Um diese Aufgabe zu lösen betrachten wir eine Stromlinie von der Oberfläche des Wassers bis hin zur Tiefe h. Beachte, dass eine Stromlinie als Tangente an die Geschwindigkeitsvektoren definiert ist. Da bei einem ruhenden Fluid die gesamten Vektoren null sind, kann man eine Stromlinie letztlich auf beliebigen Wegen ziehen. Das Bezugsniveau für die Lageenergien legen wir in die betrachtete Tiefe, in der der Druck bestimmt werden soll. Dieser Tiefe ist also die Höhe null zugeordnet und der Wasseroberfläche die Höhe h. Der statische Druck an der Wasseroberfläche entspricht dem Umgebungsdruck p$_u$. Damit sind folgende Werte bekannt:


\begin{table}[h]
\centering
\begin{tabular}{l|l|l}
& Zustand 1 (Wasseroberfläche) & Zustand 2 (Tiefe) \\ \hline
Höhe & $h_1 = h$ & $h_2 = 0$ \\ \hline
Geschwindigkeit & $v_1 = 0$ & $v_2 = 0$ \\ \hline
Statischer Druck & $p_1 = p_u $ & $p_2 =$ gesucht \\ \hline
\end{tabular}
\end{table}



\ifshowme
\antwort{
\begin{align*}
\underbrace{p_1}_{p_u} + \frac{1}{2} \rho \underbrace{v_1^2}_{=0} + \rho g \underbrace{h_1}_{=h} &= p_2 + \frac{1}{2} \rho \underbrace{v_2^2}_{=0} + \rho g \underbrace{h_2}_{=0} \\
p_2 &= p_u + \rho g h
\end{align*}
}
\else
\fi

\pagebreak




\subsection*{Aufgabe 2}
Der Austrittsdiffusor einer Turbine erweitere sich von der Eintrittsfläche $A_1$ auf die Austrittsfläche $A_2$. Der Diffusoreintritt liege um die Höhe $h$ über dem Wasserspiegel.
Für die Werte $h = 2 ~\text{m}$, $A_1 = 5 ~\text{m}^2$, $A_2 = 25 ~\text{m}^2$, $\rho = 1 \cdot 10^3 ~\text{kg/m}^3$, $p_0 = 1 ~\text{bar}$ bestimme man den maximalen Volumenstrom $\dot{V}$, sodaß der Druck $p_1$ im Diffusoreintritt
den Dampfdruck $p_D \approx 0$ bar des Wassers gerade nicht unterschreitet.
\begin{figure}[h]
\centering
\resizebox{.65\textwidth}{!}{
\begin{tikzpicture}[>=stealth, scale=1.2, font=\small]

    % Parametrisierung der Koordinaten für einfache Anpassungen
    \def\rOne{1}       % Radius oben (A1)
    \def\rTwo{2.5}     % Radius unten (A2)
    \def\hTop{3.5}     % Höhe der abgeschnittenen Rohrleitung
    \def\hEntry{2}     % Höhe des Diffusoreintritts
    \def\hWater{-1}    % Höhe des Wasserspiegels
    \def\hExit{-2.5}   % Höhe des Diffusoraustritts
    \def\dimX{-3.2}    % X-Position der Bemaßungslinien

    % Wasserspiegel (Trennlinie Luft / Wasser)
    \draw (-4.5, \hWater) -- (4.5, \hWater);
    
    % Beschriftung Luft & Wasser
    \node[above right] at (-4.5, \hWater) {\large Luft};
    \node[below right] at (-4.5, \hWater) {\large Wasser};

    % Symbol für Wasseroberfläche (kleines Dreieck)
    \draw (3.2, \hWater) -- (3.4, \hWater+0.25) -- (3.0, \hWater+0.25) -- cycle;

    % Symmetrieachse (strichpunktiert)
    \draw[dotted,line width=1pt] (0, \hTop+0.5) -- (0, \hExit-0.5);

    % --- Diffusor Kontur ---
    
    % Abgeschnittenes Rohr oben (Wellenlinie)
    \draw (-\rOne, \hTop) sin (-0.5, \hTop+0.1) cos (0, \hTop) sin (0.5, \hTop-0.1) cos (\rOne, \hTop);
    
    % Vertikales Rohrstück oben
    \draw (-\rOne, \hTop) -- (-\rOne, \hEntry);
    \draw (\rOne, \hTop) -- (\rOne, \hEntry);
    
    % Eintrittsfläche A1 (Trennlinie)
    \draw (-\rOne, \hEntry) -- (\rOne, \hEntry);
    
    % Konischer Teil des Diffusors
    \draw (-\rOne, \hEntry) -- (-\rTwo, \hWater);
    \draw (\rOne, \hEntry) -- (\rTwo, \hWater);
    
    % Zylindrischer Teil unten (unter Wasser)
    \draw (-\rTwo, \hWater) -- (-\rTwo, \hExit);
    \draw (\rTwo, \hWater) -- (\rTwo, \hExit);
    
    % Austrittsfläche A2 (Bodenlinie)
    \draw (-\rTwo, \hExit) -- (\rTwo, \hExit);


    % --- Beschriftungen ---
    \node[scale=1.3] at (-0.5*\rOne, \hEntry+0.35) {$p_1$};
    \node[scale=1.3] at (-0.5*\rOne, \hEntry-0.35) {$A_1$};
    \node[scale=1.3] at (\rTwo+0.5, 0.5) {$p_0$};
    \node[scale=1.3] at (-0.5*\rTwo, \hExit+0.35) {$A_2$};

    % Pfeil für den Volumenstrom (V-Punkt)
    \draw[->,very thick] (0, \hEntry-0.5) -- (0, \hWater+0.2) node[midway, right=2pt] {$\dot{V}$};


    % --- Bemaßungen ---
    
    % Hilfslinien für Bemaßung
    \draw (-\rOne-0.2, \hEntry) -- (\dimX-0.2, \hEntry);
    \draw (-\rTwo-0.2, \hExit) -- (\dimX-0.2, \hExit);

    % Bemaßung h
    \draw[<->,very thick] (\dimX, \hEntry) -- (\dimX, \hWater) node[midway, left=2pt] {\large $h$};
    
    % Bemaßung D
    \draw[<->,very thick] (\dimX, \hWater) -- (\dimX, \hExit) node[midway, left=2pt] {\large $D$};

\end{tikzpicture}
}
\caption{Abbildung eines Austrittsdiffusors}
\end{figure}

\ifshowme
\antwort{
\begin{align*}
p_1 + \rho g h_1 + \frac{1}{2}\rho v_1^2 &= p_2 + \rho g h_2 + \frac{1}{2}\rho v_2^2 \\
p_1 + \rho g h_1 + \frac{1}{2}\rho \left( \frac{\dot{V}}{A_1} \right)^2 &= p_2 + \rho g h_2 + \frac{1}{2}\rho \left( \frac{\dot{V}}{A_2} \right)^2 \\
p_1 + \rho g h_1 + \frac{1}{2}\rho \left(\frac{\dot{V}}{A_1}\right)^2 &= p_0  + \rho g D - \rho g D + \frac{1}{2}\rho \left( \frac{\dot{V}}{A_2} \right)^2 \\
p_1 + \rho g h_1 + \frac{1}{2}\rho \left( \frac{\dot{V}}{A_1} \right)^2 &= p_0  + \frac{1}{2}\rho\left(  \frac{\dot{V}}{A_2} \right)^2 \\
\end{align*}
Wir setzen $p_1$ = 0
\begin{align*}
0 + \rho g h_1 + \frac{1}{2}\rho \left( \frac{\dot{V}}{A_1} \right)^2 &= p_0  + \frac{1}{2}\rho \left( \frac{\dot{V}}{A_2} \right)^2 \\
\end{align*} 
Nach Volumenstrom umstellen und einsetzen:
\begin{align*}
\frac{1}{2}\rho \left( \frac{\dot{V}}{A_1} \right)^2 - \frac{1}{2}\rho \left( \frac{\dot{V}}{A_2} \right)^2 &= p_0 - \rho g h\\
\left( \frac{\dot{V}}{A_1} \right)^2 - \left( \frac{\dot{V}}{A_2} \right)^2 &= \frac{p_0 - \rho g h}{\frac{1}{2}\rho}\\
\dot{V} &= \sqrt{\frac{p_0 - \rho g h}{\frac{1}{2} \rho \cdot\left[ \left(\frac{1}{A_1} \right)^2 - \left(\frac{1}{A_2} \right)^2 \right] }} \\
\dot{V} &= \sqrt{\frac{100.000 - 1.000 \cdot 9,81 \cdot 2}{\frac{1}{2} \cdot 1.000 \cdot\left[ \left(\frac{1}{5~\text{m}^2} \right)^2 - \left(\frac{1}{25~\text{m}^2} \right)^2 \right] }}
\end{align*}
\begin{align*}
\dot{V} \geq 64,7 ~\text{m}^3\text{/s}
\end{align*}
}
\else
\fi


\subsection*{Aufgabe 3}
In einer Gasleitung ($\rho_{Gas} = 1, 18 ~\text{kg/m}^3$
, $\nu_{Gas} = 15 \cdot 10^{-6} ~\text{m}^2\text{/s}$) von 0,2 km Länge und
150 mm Durchmesser strömt das Gas mit einer Geschwindigkeit von 20 m/s. Die Rauhigkeit
$k$ sei 0, 075 mm (geschweißte Stahlrohre).
Man berechne den Druckabfall $\Delta p$. Die Strömung kann als inkompressibel angenommen
werden.

\ifshowme
\antwort{
Reynoldszahl:
\begin{align*}
Re &= \frac{v \cdot l}{\nu} \\
 &= \frac{20~\text{m/s} \cdot 0,15 m}{15 \cdot 10^{-6} ~\text{m$^2$/s}} \\
 &= 2 \cdot 10^5 \rightarrow ~\text{turbulent}
\end{align*}
Relative Wandrauhigkeit:
\begin{align*}
\frac{d}{k} = \frac{150}{0,075} = 2.000
\end{align*}
$\xi_{Rohr}$ aus Moody-Diagramm ablesen \RA 0,02


\begin{align*}
\Delta p &= \xi_{Rohr} \cdot \frac{l}{d} \cdot \frac{1}{2}\rho \overline{v}^2 \\
 &= 0,02 \cdot \frac{200}{0,15} \cdot \frac{1}{2} \cdot 1,18 \cdot 20^2 \\
 &= 6.293,3 ~\text{Pa}
% &= 5.82 \cdot 10^{-2}~\text{bar}
\end{align*}
}
\else
\fi

\subsection*{Aufgabe 4}
Was versteht man unter der Nusselt-Zahl und der Prandtl-Zahl? Versuchen Sie die Begriffe in anschaulichen Beispielen darzustellen

\ifshowme
\antwort{

Die \textbf{Nußelt-Zahl} drückt aus, um welchen Faktor die Wärmeübertragung aufgrund der Konvektion stärker ist als wenn reine Wärmeleitung wirken würde.
\begin{align*}
\text{Nu} = \frac{\text{Wärmetransport Konvektion}}{\text{Wärmetransport Leitung}} = \frac{\alpha \cdot L}{\lambda}
\end{align*}

\pp
Die \textbf{Prandtl-Zahl} kann wie folgt interpretiert werden:

\begin{itemize}
\item Eine hohe Prandtl-Zahl bedeutet, dass das Fluid eine hohe Viskosität und eine niedrige Temperaturleitfähigkeit hat.
\item Eine niedrige Prandtl-Zahl bedeutet, dass das Fluid eine niedrige Viskosität und eine hohe Temperaturleitfähigkeit hat.
\end{itemize}

\par \vspace*{-4mm}
\begin{align*}
\text{Pr} = \frac{\text{Impulstransport}}{\text{Wärmetransport}} = \frac{\nu}{\alpha}
\end{align*}

Da der Impulstransport durch das Geschwindigkeitsfeld, der Wärmetransport durch das Temperaturfeld bestimmt ist, verbindet die Prandtl-Zahl die beiden für den Wärmeübergang maßgebenden Felder. Die Prandtl-Zahl ist somit ein Maß für das Verhältnis der Dicken von Strömungsgrenzschicht zu Temperaturgrenzschicht.

}
\else
\fi



\end{document}
